\section{A lower bound from the price of a shared label}\label{sec:boundary}

Consider a square matrix product of side \(n\), with its two inputs
and its output at the tips of a tripod with three unit arms. Each of
the three matrix indices is shared by two tips at distance \(2\), so
the product tensor has price \(6\log n\). A rank-\(r\) decomposition
supplies that tensor from a label with \(r\) values shared by all
three tips. If the rate of this shared label is \(8/3\), its price is
\((8/3)\log r\), and price monotonicity gives \(r\ge n^{9/4}\).
We will find such a tripod in the family constructed in
Section~\ref{sec:motion}. We first test the shared-label rate when one
arm is much longer than the other two. The following principle carries
the resulting equality to three equal arms.

\begin{proposition}[Tripod principle]\label{prop:tripod}
Over an algebraically closed field, place three sites at pairwise
distances \(a_i+a_j\) for \(i\ne j\), where \(a_1,a_2,a_3\ge0\).
Their common-label rate \(D(a)\) satisfies
\[
       D(a)\ \ge\ D_*(a)
       :=a_1+a_2+a_3-\bigl(a_1^{-1}+a_2^{-1}+a_3^{-1}\bigr)^{-1}.
\]
The reciprocal term is zero when any \(a_i=0\).

Now fix three infinite rays in a spectral space that form a metric
tripod with origin \(O\), and suppose that their prices are
homogeneous under radial dilation. In \(D(a)\), place site \(i\) on
ray \(i\) at distance \(a_i\) from \(O\). If \(D(a)-D_*(a)\)
vanishes at one positive triple, it vanishes at every positive
triple. The same conclusion holds if the defect tends to zero as one
arm tends to infinity with the other two held at fixed positive lengths.
\end{proposition}

The first assertion applies to any placement with the specified
distances. The second compares placements on a single tripod.
Sections~\ref{sec:profiles}--\ref{sec:contact} prove both assertions by
studying how the price of polynomial multiplication changes with its
degree and with the positions of its coefficients.

\subsection{A decoder for rectangular multiplication}

Multiply an \(s\times q\) matrix by a \(q\times s\) matrix over the
algebraically closed field \(F\) of Section~\ref{sec:motion}, taking
\(s=N^{1/2}\) and \(q=N^{1/4}\) for a fourth power \(N\). The product
has \(N\) entries. Place its output at \(A\) and its inputs at
\(C,D'\), with respective arm lengths \(c,a,b\). The outer indices
give wires of size \(s\) on \(AC,AD'\), and the inner index gives a wire
of size \(q\) on \(CD'\). The price of the product, divided by
\(\log N\), therefore satisfies
\[
   c+\tfrac34(a+b)
   \ \le\ c+a+b-\frac{ab}{a+b}
   \ \le\ D_*(c,a,b)\ \le\ D(c,a,b).
\]
Here \(ab/(a+b)=0\) if \(a=b=0\). This gives the rectangular bound
associated with half-width multiplication \cite{uznanski,abvz}.

In particular, for every \(\eta>0\) there are arbitrarily large
such \(N\) for which the exact tensor rank is at most \(N^{1+\eta}\).
Indeed, start with a sufficiently large fourth power \(N_0\)
and a full diagonal of size \(\lceil N_0^{1+\eta/2}\rceil\).
The displayed inequality at every spectral point gives strict
spectral comparison with this diagonal. Each tensor power of the product
is a restriction of a bounded number of copies of the corresponding
power of the diagonal. At sufficiently large powers the exponent
\(1+\eta\) absorbs both this multiplicity and the rounding.

\begin{lemma}[A decoder with short rows]\label{lem:secondary}
For every \(\epsilon>0\) there is a \(\delta>0\) such that for
unbounded fourth powers \(N\) the \(N\) output bilinear forms \(B_e\)
of this product admit decompositions
\[
    B_e=S_e+R_e,\qquad
    \operatorname{rank}S_e\le N^\epsilon,\qquad
    \dim\langle R_e:1\le e\le N\rangle\le N^{1-\delta}.
\]
\end{lemma}

The two bounds have different uses. A short sum of simple products
computes each \(S_e\) once the output index \(e\) is shared.
The corrections \(R_e\) may be expensive individually, but their
entire span can be sent across a single wire. To obtain the bounds
uniformly in the base format, we implement a tensor power of a decoder
using one work register. For related nonrigidity estimates for
Kronecker products, see \cite{almanrigidity,kivva}.

\begin{proof}
It suffices to take \(0<\epsilon<1\). Fix an integer \(K>10/\epsilon\).
For a base format with \(d\) output entries and rank
\(d\le r\le d^{1+\eta}\), let \(M\) be the \(d\times r\) matrix that
expresses its outputs in terms of \(r\) simple bilinear products.
We will choose \(\eta\) below. The \(Km\)-th tensor power has
\(N=d^{Km}\) entries. Pad each decoder block \(M^{\otimes m}\) with
zero rows and columns to side \(2^t\), where \(t=\lceil m\log_2 r\rceil\).
The full padded decoder has side \(T=2^{Kt}\). An index consists of
\(K\) registers of \(t\) bits each.

Apply the first block matrix to the first register. For each subsequent
block, swap its register with the first, apply the block matrix there,
and swap back. An exchange of registers performs \(t\) simultaneous
exchanges of pairs of bits. A bit exchange transposes two of four
letters, \(01\) and \(10\). Its matrix factors into a fixed number of
elementary shears and diagonal signs, over any field. Thus a register
exchange consists of a fixed number of stages, each applying the same
local shear or diagonal sign at every bit position.

For a shear, let the high letter be the target of its off-diagonal
entry. Immediately before and after the stage, project onto the
coordinates whose register pair contains between \(t/4-\rho t\) and
\(t/4+\rho t\) high letters, for a fixed small \(\rho>0\).
The discarded fraction of coordinates is at most \(e^{-c_\rho t}\)
for large \(t\), by the binomial tail. Telescoping over the \(O(K)\)
projections bounds the rank of the discarded part of the decoder by
\(O(K)Te^{-c_\rho t}\le T^{1-\gamma}\), with \(\gamma=\gamma(K,\rho)>0\).
On a path retained at a shear stage, every change creates a high
letter, so at most \(2\rho t\) positions change.

Only the first register is acted on by a block matrix. Every other
register participates in just two exchanges, so its input lies within
Hamming distance \(C\rho t\) of its output on a retained path, for an
absolute constant \(C\). Each output row of the retained decoder
therefore has at most
\[
       2^t\left(\sum_{j\le C\rho t}\binom{t}{j}\right)^{K-1}
       \ \le\ T^{1/K+o_{\rho\downarrow0}(1)}
\]
nonzero entries. Choose \(\rho\) so that this exponent is below
\(\epsilon/2\), then choose \(\eta\) sufficiently small. Since
\(T\le 2^K N^{1+\eta}\), for all large \(m\) the two bounds become
\(N^\epsilon\) entries per retained row and rank \(N^{1-\delta}\)
for the discarded matrix, with \(\delta=\gamma/2\). Remove the
zero padding and apply both matrices to the list of simple bilinear
products. The retained rows give the forms \(S_e\), and the image of
the discarded row space spans the \(R_e\).
\end{proof}

\subsection{Two copies of the output near a collision}

Choose the collision in Section~\ref{sec:motion} where site \(4\)
approaches site \(3\), and denote sites \(1,2,3,4\) by \(A,B,C,D'\).
At parameter \(r\) write \(L=2\ell\) and \(h=L-r\), with \(0<h<L\).
The distances within \(AB\) and \(CD'\) are both \(h\), and all four
cross distances are \(L\). Write \(D_\ell(h)\) for the four-site
common-label rate. At the collision \(D_\ell(0)=L\), and at the
generic four-site placement \(D_\ell(L)=D_\ell\). Convexity gives
\[
                   D_\ell(h)\le L+\frac hL(D_\ell-L).
\]

Duplicate the output index \(e=(i,i')\) of the rectangular product at
\(A\) and \(B\), leaving its two inputs at \(C,D'\). The resulting tensor is
\[
 U_N=\sum_e e_e\otimes e_e\otimes B_e
     =\sum_{\substack{1\le i,i'\le s\\1\le j\le q}}
        e_{ii'}\otimes e_{ii'}\otimes e_{ij}\otimes e_{ji'} .
\]
The index \(i\) is a shared label at \(A,B,C\), the index \(i'\) is a
shared label at \(A,B,D'\), and \(j\) travels on the short wire \(CD'\).
Writing \(D_{ABC},D_{ABD'}\) for the two triple rates, we have
\[
   \frac{H(U_N)}{\log N}=\frac{D_{ABC}+D_{ABD'}}2+\frac h4 .
\]
Both triples have arms \((h/2,h/2,L-h/2)\). Their reciprocal
correction in Proposition~\ref{prop:tripod} is at most \(h/4\), so
both rates are at least \(L+h/4\).

Use the decomposition \(B_e=S_e+R_e\) from Lemma~\ref{lem:secondary}.
In units of \(\log N\), the tensor formed from the \(S_e\) costs at
most \(D_\ell(h)+\epsilon h\). Share the label \(e\) at all four sites,
then use a wire of size at most \(N^\epsilon\) on \(CD'\) to supply
the simple products for the specified \(e\). The tensor formed from
the \(R_e\) has matrix rank at most \(N^{1-\delta}\) across the cut
\(AB\mid CD'\). A wire between \(A\) and \(C\) supplies this rank, and
wires within the two pairs deliver the legs at \(B\) and \(D'\), each
of dimension at most \(N\). This costs at most \((1-\delta)L+2h\).

Take \(h/L<\delta/2\). The correction then costs less than \(L\), and
\(D_\ell(h)\ge L\) by restriction to a cross pair. For a fixed placement
the ordinary-sum price bound costs an additive constant independent of
\(N\). Sending \(N\) to infinity along the formats of Lemma~\ref{lem:secondary} and
then using Equation~\eqref{eq:motion-upper} at fixed \(h/L\) gives
\[
  \frac{D_{ABC}+D_{ABD'}}2+\frac h4
  \ \le\ D_\ell(h)+\epsilon h
  \ \le\ L+\bigl(\tfrac12+\epsilon+o_{\ell\downarrow0}(1)\bigr)h .
\]
Subtract the lower bound \(D_{ABD'}\ge L+h/4\) to obtain
\[
        D_{ABC}\le L+h/4+\bigl(2\epsilon+o_{\ell\downarrow0}(1)\bigr)h.
\]

The first three sites lie on the fixed rays \(z_1,z_2,z_3\) out of
the depth-zero point. Take the tangent tripod there by placing sites
at depths \(t a_i\), dividing prices by \(t\), and letting \(t\downarrow0\)
in a spectral ultralimit. On these fixed rays the rescaled price of
every nonzero tensor has a limit, since its original price is convex
in \(t\) and vanishes at zero. The limiting prices are homogeneous
under radial dilation.

Set \(t=h/2\) and hold \(b=2L/h-1\) fixed. The three depths become
\(t(1,1,b)\). Dividing the last inequality by \(t\) gives
\(D(1,1,b)\le b+3/2+4\epsilon\), whereas
\[
             D_*(1,1,b)=b+\tfrac32+\frac{1}{4b+2}.
\]
Thus \(0\le D(1,1,b)-D_*(1,1,b)\le4\epsilon\) for every sufficiently
large rational \(b\), and for real \(b\) by continuity. Sending \(b\)
to infinity and then \(\epsilon\) to zero gives contact at infinity
on this fixed tripod. Proposition~\ref{prop:tripod} now gives
\(D(1,1,1)=3-1/3=8/3\). The three-wire price comparison at the start
of this section yields \(r\ge n^{9/4}\), proving Theorem~\ref{thm:motion}.
